\begin{statement}[C.8.4; Proved: the first three normalized duplication steps]
Put $u=t-1$. For $n=2^k$, $k=1,2,3$, and $nP\ne O$,
\begin{equation}\label{eq:ca-ladder}
 \left(\frac{N(P)}{\sqrt2}\right)^{n^2}
 =\frac{3^{3/4}}{2n}\,
       \frac{N(nP)|u|^{-1/4}}{|A_k(u)|},
\end{equation}
where
\begin{align*}
 A_1(u)&=1,\qquad A_2(u)=1-5u/27,\\
 A_3(u)&=1-\frac{217}{27}u-\frac{20486}{729}u^2
       +\frac{10910}{19683}u^3
       +\frac{161383}{177147}u^4
       -\frac{245}{14348907}u^5.
\end{align*}
On $|u|\le1/28$ we have $|A_3(u)|\ge1-10|u|$.
For every $k\ge1$, not just these first three steps, define
\[
 E_k(t)=\frac{|\Delta_E|^{(4^k-1)/12}|t-1|^{1/4}}
 {2^{4^k/2}|\psi_{2^k}(P)|}.
\]
Then $(N(P)/\sqrt2)^{4^k}=E_k(t)N(2^kP)|t-1|^{-1/4}$,
and the removable limiting value is
\[
 \lim_{t\to1}E_k(t)=\frac{3^{3/4}}{2^{k+1}}
                 =2^{1-k}\frac{3^{3/4}}4.
\]
\end{statement}
\begin{proof}
Substitute the exact division polynomials from Statement C.4.2
in \eqref{eq:cs-division}. Equivalently, remove
$t^{2(n^2-1)/3}(t-1)^{n^2/4}$ from $\psi_n(P)$ and divide
the resulting polynomial by its value at $t=1$. For $n=2,4,8$
this gives the displayed polynomials and the constant
$3^{3/4}/(2n)$. All operations are rational polynomial operations;
\path{code/C_ladder_algebra.py} reproduces them and records
the next two steps as well. The general identity is C.4.2.
For its limit write $\mu=(1-\tau)/2+h$ as in C.3.1. Since
$n=2^k$ is even, $n\mu$ is congruent to $nh$ modulo the
period lattice. Hence
\[
 N(nP)=2\pi|\eta(i)|^2 n|h|(1+o(1)),\qquad
 t-1=\tfrac94\pi^4E_4(i)h^4(1+o(1)).
\]
The identity $E_4(i)^3=1728\eta(i)^{24}$ then gives
$N(nP)|t-1|^{-1/4}\to2n/3^{3/4}$.
Since the left side of the general ladder tends to one,
this proves the stated limit for every $k$.
The triangle inequality gives
$|A_3(u)|\ge1-|u|M$, where
\[
 M\le\frac{81}{10}+\frac{29}{28}
                +\frac1{28^2}+\frac1{28^3}+\frac1{28^4}<10.
\]
Here we bounded respectively the five absolute coefficients by
$81/10,29,1,1,1$. This is a rational inequality on one disc,
not a subdivision or a root search.
\end{proof}
